\documentclass[11pt,a4paper]{article}
\usepackage[margin=22mm]{geometry}
\usepackage{amsmath,amssymb,booktabs,array}
\usepackage[hidelinks]{hyperref}
\setlength{\parindent}{0pt}
\setlength{\parskip}{5pt}
\raggedright
\emergencystretch=3em
\begin{document}
\begin{center}
{\Large T.A amount allowances and a consistent allocation choice}\\[4pt]
8 October 2026 --- methodology-only analytical continuation
\end{center}
\textbf{Status.} T.A is approved. The new Q.J.F.P question concerns the time
meaning of competition proportions, not a held source rate. Both alternatives
have qualified local examples; no whole joined solver or gradient is adopted.

\section*{1. T.A state, units and conservation}
Use metres, seconds and continuous vehicle-equivalent mass (eq). Density
$\rho$ has units eq/m and flow has units eq/s. The common triangular field is
\begin{equation}
 F(\rho)=\min\{u\rho,w(\kappa-\rho)\},\quad
 u=30,\ w=5,\ \kappa=3/20,\quad \rho_c=3/140,\ Q=9/14.
\end{equation}
Sections have $L=30$ m and public interval $T=1$ s. Sampled source sending and
movement masses are
\begin{equation}
 S_a=\min\{n_a(0),QT\},\qquad d_{ab}=\theta_{ab}S_a,\quad
 \sum_b\theta_{ab}=1.
\end{equation}
The choice fractions use the approved prior-period lookahead Edie observations;
they are not newly sampled at interior wave events. For actual accepted port
flow $f_{ab}$, define the unspent numerical claim
\begin{equation}
 b_{ab}(t)=d_{ab}-\int_0^t f_{ab}(s)\,ds\ge0.
\end{equation}
This is request accounting, not a second traffic population or assigned vehicle
identity. Pair each transfer once. When a claim reaches zero, its movement
cannot discharge further this interval. There is no imposed bound
$f_{ab}\le d_{ab}/T$. For constant local rate $f>0$, the remaining amount
is used after $b/f$ seconds; stop at that exact event rather than an arbitrary
threshold. Claims refresh at the public clock, without rerouting unspent mass.

The donor's total outflow uses its own physical field. All target inputs spend
one receiving resource, and outgoing debits must not spend donor mass twice.
Under the existing upstream-port geometry and $L=uT$, post-clock incoming
traffic cannot traverse a whole section within the interval. Consequently
\begin{equation}
 \sum_b\int_0^T f_{ab}\,dt\le\min\{n_a(0),QT\}=S_a.
\end{equation}
Individual $d_{ab}$ do not follow from that total bound. They need the claim
ledger. E1 accumulates completed actual transfers and creates zero-population
effects at the public clock, never at interior claim exhaustion.

\newpage
\section*{2. Exact allowance exhaustion and stop waves}
One donor starts with jam density on $[86/3,30]$ and vacuum upstream; its
mass is $1/5$. Upstream admission is closed. Desired lateral and through
masses are $1/20$ and $3/20$. Use sampled lateral choice fraction $1/4$;
the through target has zero receiving, while the lateral target is free.
Before the claim is fulfilled, the source-oriented local rate is
\begin{equation}
 \beta=Q/4=9/56,\quad \rho_\beta=\kappa-\beta/w=33/280,\quad
 c_\beta=\beta/\rho_\beta=15/11.
\end{equation}
It is a declared source-orientation regime, not a unique general allocator
selected by T.A. The outgoing queue trace has sending $Q$. The initial
jam-to-$\rho_\beta$ release wave travels upstream at $-5$, reaching the
rear at
\begin{equation}
 t_0=(4/3)/5=4/15.
\end{equation}
Before $t_0$, vacuum ends at $86/3$, then jam extends to $30-5t$, followed
by $\rho_\beta$ to the exit. Afterwards the vacuum-to-$\rho_\beta$ rear
travels at $c_\beta$. The lateral claim is fulfilled at
\begin{equation}
 t_q=(1/20)/\beta=14/45,\quad
 x_q=86/3+c_\beta(t_q-t_0)=316/11.
\end{equation}
At $t_q$, stop lateral flow. Through receiving is still zero, so donor outflow
becomes zero without reassigning its remaining through claim. A jam wave
travels backward from 30 at $-5$, meeting the moving rear after
\begin{equation}
 t_c-t_q=\frac{30-x_q}{5+c_\beta}=1/5,\qquad
 t_c=23/45,\quad x_c=29.
\end{equation}
Thereafter the vacuum-to-jam rear is stationary. At $T=1$ the donor is
jammed on $[29,30]$, with mass $3/20$. Exactly $1/20$ transferred.
At service start and stop use the post-event one-sided port trace: its new
wave initially has zero length, so this changes no integrated field mass.
In the initially empty recipient, density is $\beta/u=3/560$ on
$[u(1-t_q),u]=[62/3,30]$ at the completed clock. Its mass is $1/20$;
no positive-measure traffic has exited it before that clock. Donor/recipient
counts therefore pair exactly.

Without the T.A amount check, repeated source-oriented rate $\beta$ could
continue for the entire second: remaining donor mass $1/5-\beta=11/280>0$.
The ordinary rear then ends at $89/3$, with $\rho_\beta$ on its last $1/3$ m.
Accepted lateral mass would be $9/56>1/20$, although aggregate traffic is still
conserved. This refutes a naive repeated-demand transplantation, not native
finite IT or continuous mass. This example resolves concrete wave interactions
at depletion; arbitrary depleted multi-port modes remain unqualified.

\newpage
\section*{3. Actual ramp/two-lane geometry for the allocation comparison}
Donor 1 advances through to shoulder recipient A. Donor 2 advances laterally
to A with sampled fraction $\theta=1/4$, or through to free recipient B with
fraction $3/4$. A finite KW ramp also enters A. There is no same-section
lateral transfer followed by another longitudinal advance.

Initially donor 1 is uniformly free at $\rho_1=1/150$, with sending
$D_1=1/5$; donor 2 and the ramp start critical, sending $D_2=D_r=Q$.
Compatible prescribed inlet forcing is $1/5,Q,Q$, respectively. These are
physical boundary inputs; the ramp admission can be realized by M.A under
adequate offered flow and permission. No downstream wave reaches its upstream
meter in this one-second case. Use supplied prior Edie observations such
that $\theta=1/4$ (e.g. 30 and 15 m/s with relaxation time 2 s). Their reached
control history is not claimed.

Recipient A has density $\rho_A=9/100$ on $[0,1]$ and vacuum ahead. An old
zero-passing effect starts at 1, with held desired speed $s=10/3$. Own ahead
traffic is vacuum, so K.A permits exactly $s$. The queue grows with the record:
\begin{equation}
 z(t)=1+\tfrac{10}{3}t,\quad F(\rho_A)=s\rho_A=3/10.
\end{equation}
Its own inlet supply $R=3/10$ is constant. This receiving history is derived
from the same queue/record field, not an external capacity penalty.
With illustrative fixed ramp share $p=1/3$, both groups request more than
their shares. N.A awards
\begin{equation}
 f_r=1/10,\qquad R_m=1/5\quad\text{for all non-ramp inputs together}.
\end{equation}
Recipient B starts free at density $(27/56)/30=9/560$, with free exit; its
inflow/outflow is $27/56$. All allowances are as sampled:
\begin{equation}
 d_{1A}=1/5,\quad d_{2A}=9/56,\quad d_{2B}=27/56,\quad d_{rA}=9/14.
\end{equation}
No movement is fulfilled before the interior times considered below; $d_{2B}$
is exactly fulfilled at the ending clock. Thus the allocation distinction
exists even before a quota-depletion policy is needed.

\textbf{P.A.} Retain sampled non-ramp proportions $1/5:9/56=56:45$ as
competition weights. Its accepted rates are
\begin{equation}
 f_{1A}^{A}=56/505,\quad f_{2A}^{A}=9/101,\quad f_{2B}^{A}=27/56.
\end{equation}
\textbf{P.B.} Solve the source-oriented wave junction consistently. Both
donors are restricted and therefore have congested boundary sending $Q$.
Its oriented proportions are $Q:(Q/4)=4:1$, giving
\begin{equation}
 f_{1A}^{B}=4/25,\quad f_{2A}^{B}=1/25,\quad f_{2B}^{B}=27/56.
\end{equation}
Both give the same $R_m$ and respect all amount ceilings.

\newpage
\section*{4. Compatible own fields, counts and births for both allocations}
For candidate $k\in\{A,B\}$, donor 1's queue trace is
$\rho_{1k}^*=\kappa-f_{1A}^{k}/w$. Its shock from the initial free state has
speed
\begin{equation}
 \sigma_{1k}=\frac{f_{1A}^{k}-1/5}{\rho_{1k}^*-1/150}<0.
\end{equation}
It lies between $-w$ and $u$. At $t=1$ donor 1 has its initial density on
$[0,30+\sigma_{1k}]$ and $\rho_{1k}^*$ to its port. The wave never reaches
the inlet. Its own sending is then $Q$; this is distinct from actual service.
Donor 2's total outflow is $g_{2k}=f_{2A}^k+27/56<Q$, with queue density
$\rho_{2k}^*=\kappa-g_{2k}/w>\rho_c$. Its shock travels at $-5$, leaving
initial critical density on $[0,25]$ and $\rho_{2k}^*$ on $[25,30]$.
The ramp similarly has queue density $\kappa-(1/10)/w=13/100$ on $[25,30]$.

Direct field integration equals paired conservation in every lane:
\begin{center}
\begin{tabular}{lcc}
\toprule
Ending mass & P.A & P.B\\
\midrule
Donor 1 & $146/505$ & $6/25$\\
Donor 2 & $4041/5656$ & $1069/1400$\\
Ramp & $83/70$ & $83/70$\\
Recipient A & $39/100$ & $39/100$\\
Recipient B & $27/56$ & $27/56$\\
\bottomrule
\end{tabular}
\end{center}
The two donor-mainline masses sum to $281/280$ under both. Initial mass plus
inlets minus external exits equals the total ending inventory; internal
transfers cancel. Descriptors add no population.

For recipient A, local count gauge is
\begin{equation}
 C_A(t)=9/100+(3/10)t,\quad
 N_A(t,x)=C_A(t)-\int_0^x\rho_A(t,y)dy.
\end{equation}
The old record has $N_A(t,z(t))=0$ throughout. At $t=1$, $z=13/3$, the queue
mass and gauge both equal $39/100$. Thus both own-field motion/count histories
are consistent, although entry composition differs.

For an illustrative eligible lateral counter $H_0=47/50$ and threshold one,
\begin{equation}
 H_A(1)=47/50+9/101=5197/5050>1,\qquad
 H_B(1)=47/50+1/25=49/50<1.
\end{equation}
P.A creates one new E1 effect, attached at recipient A's completed inlet count
$39/100$; P.B creates none. The distinction is not a same-law hidden-information
witness: these are different temporal allocation completions.

\newpage
\section*{5. A local P.B wave-junction completion beyond naive substitution}
Derive a one-fork/two-recipient kernel with $0<\theta<1$, two positive input
demands $D_i\le Q_i$, recipient B unconstrained, receiving $R_m$ at A,
and pending amount claims. This is a reviewer derivation from source IT and
ordinary boundary-wave admissibility, not a generic FIFO theorem.
Here ``B unconstrained'' means its supply is at least $(1-\theta)Q_2$;
this holds for the common-capacity free recipient in the example.

Let $\gamma\in[0,1]$ be A's common IT receiving factor and $\xi_i$ auxiliary
interior demands. Input 1 sends $\gamma\xi_1$. Input 2 sends
$\theta\gamma\xi_2$ to A and $(1-\theta)\xi_2$ to B. A restricted incoming
road's stationary boundary state is congested: its interior sending is capacity.
An unrestricted incoming road retains total accepted flow equal to $D_i$;
an admissible interior state can have demand between $D_i$ and $Q_i$.
Therefore
\begin{align}
 f_{1A}(\gamma)&=\min\{D_1,\gamma Q_1\},\\
 \xi_2(\gamma)&=\min\left\{Q_2,\frac{D_2}{1-\theta+\theta\gamma}\right\},\\
 f_{2A}(\gamma)&=\theta\gamma\xi_2(\gamma),\qquad
 f_{2B}(\gamma)=(1-\theta)\xi_2(\gamma).
\end{align}
Define $H(\gamma)=f_{1A}(\gamma)+f_{2A}(\gamma)$. If
$R_m\ge H(1)=D_1+\theta D_2$, use $\gamma=1$. Otherwise $H(0)=0$ and
$H$ is continuous and strictly increasing, so solve $H(\gamma)=R_m$ uniquely.
The second term has positive derivative: $\theta Q_2$ on its capacity branch,
or $\theta D_2(1-\theta)/(1-\theta+\theta\gamma)^2$ on its demand branch.
The first term is nondecreasing. All input totals are at most their incoming
demands, and all output flows obey receiving.

At a restricted input, total accepted flow is below $D_i$, so its capacity
branch is already selected. Replacing that input's demand by capacity after
its queue forms leaves the same selected flow. This is the local consistency
that a direct initial point-proportional substitution lacked.

For the actual comparison $D_1=1/5$, $D_2=Q=9/14$, $\theta=1/4$,
$R_m=1/5$, both capacity branches give
\begin{equation}
 \gamma=\frac{R_m}{Q(1+\theta)}=56/225,\quad
 f_{1A}=4/25<1/5,\quad f_{2A}=1/25,\quad f_{2B}=27/56.
\end{equation}
Input 2's total $731/1400$ is below $Q$, consistent with its queue. The source
rule holds on the congested interior states. Input 1's actual traffic field
receives the admissible shock in section 4. No externally prescribed queue
or altered ramp priority is needed to select these local flows.

\textbf{Limits.} Two constrained recipients, more lanes/inputs, zero or
unit choice fractions, depleted claims and arbitrary restriction/node contacts
require separate completion. This algebra is not an all-network theorem or
a convergence proof for the finite source step. Primary continuum general-
junction theorems cited here include FIFO coupling; they are not imported.

\newpage
\section*{6. P.A's local fixed-proportion completion and the pending choice}
For the same one-fork geometry, let sampled weights $\mu_1,\mu_2>0$ remain
constant within the public interval and $\alpha=\mu_1/(\mu_1+\mu_2)$.
The same unconstrained-B supply condition applies.
An explicitly declared local fixed-priority completion uses potential target-A
demands $D_1$ and $\theta D_2$. With unused shares transferred, define
\begin{align}
 f_{1A}&=\min\{D_1,\max(\alpha R_m,R_m-\theta D_2)\},\\
 f_{2A}&=\min\{\theta D_2,\max((1-\alpha)R_m,R_m-D_1)\},\\
 \xi_2&=\min\left\{Q_2,\frac{D_2-f_{2A}}{1-\theta}\right\},\quad
 f_{2B}=(1-\theta)\xi_2.
\end{align}
Since $f_{2A}\le\theta D_2$, $\xi_2\ge D_2$, and the donor total
$f_{2A}+f_{2B}$ is at most $D_2$. If that total is strictly below $D_2$,
then $\xi_2=Q_2$ and the donor queue is admissible. If equal, an interior
demand $\xi_2\in[D_2,Q_2]$ is admissible. Thus this kernel also has compatible
local wave states, without introducing a held $d/T$ rate ceiling.

In the comparison, $\mu_1=1/5$, $\mu_2=9/56$, $\alpha=56/101$;
both shares are below their potential demands, giving $56/505$ and $9/101$.
The second donor is capacity-restricted in total, so $\xi_2=Q$ and its
parallel through flow is $27/56$, as in the exact field table.

In this local geometry, the fixed-weight kernel is invariant to a restricted
input's demand increase. For input 1 this is the usual fixed-priority merge
property. For input 2, strict total restriction implies
$f_{2A}<D_2-(1-\theta)Q_2\le\theta D_2$; its target-A demand cap is already
nonbinding in the allocation. Increasing $D_2$ to $Q_2$ leaves its assigned
share unchanged, while the same parallel capacity branch is retained.
Unrestricted inputs do not require a queue-induced demand replacement.

This is a declared source-time priority adaptation, not a proof that native
IT already specifies it for nonuniform physical-time fields. The primary
node literature supports fixed-priority allocation; it does not choose our
sampled-demand weights or qualify every joined/depleted mode.

\textbf{Choice.} T.A determines each movement's integrated ceiling. It does
not choose whether its original request determines a standing service weight
(P.A), or source-proportional wave-junction demands determine its share (P.B).
Both local kernels can be causal, conservative and quota-compatible, yet the
flows and next-clock effects differ. Q.J.F.P therefore reopens only the
temporal meaning of IT competition, including N.A's non-ramp subgroup.
Fixed ramp priority/shared receiving, M.A/G1/E1/K.A/T.A stay approved.

P.A is provisionally preferred for the sampled task. After selection,
qualify shared-donor/multi-target allocation and depleted modes, then the
remaining motion/contact/crossing/E1 completion. No full Q, clearance, full
TST/control response, useful gradient or Stage 4 is claimed.

\section*{Primary reading and attribution}
Laval and Daganzo (2006), equations 7--13: sampled oriented amounts and IT.
Wright et al., arXiv:1601.01054, sections 2.1/3.1: demand, priority and
invariance; arXiv:1608.07623: cumulative node service/completion events, whose
snapshot execution clock is not automatically our physical clock.
Jin, arXiv:1204.6727, section 3.1: admissible stationary/interior states;
its general theorem has a different FIFO scope. Sections 2--6 above contain
reviewer derivations under their stated conditions.

\href{https://doi.org/10.1016/j.trb.2005.04.003}{Source IT}\quad
\href{https://arxiv.org/pdf/1601.01054}{Priority/node manuscript}\quad
\href{https://arxiv.org/pdf/1608.07623}{Node dynamic-system manuscript}\quad
\href{https://arxiv.org/pdf/1204.6727}{Wave-junction manuscript}
\end{document}
